\section{引言：为什么 Planning 是 Coding Agent 的核心能力}

本讲继续讨论 Codex / Claude Code 这一类 coding agent 的工具设计。讲者指出，前面介绍 Codex 的最小工具集时，容易漏掉一个看似简单但非常关键的工具：\texttt{update\_plan}。在实际开发任务里，一个 agent 面对的不是单步问答，而是一个不断展开的代码库任务：先理解 codebase，再形成计划，然后执行、检查、调整，直到任务完成。

\begin{figure}[H]
\centering
\includegraphics[width=\textwidth]{frames/frame-00-00-50.jpg}
\caption{Modern Agent 系列中，Planning Agent 与 Codex 工具设计形成一条连续线索。\protect\footnotemark}
\end{figure}
\footnotetext{视频画面时间区间：00:00:45--00:00:55。}

\begin{importantbox}{Codex 的典型任务流}
在讲者的抽象里，Codex 处理一个开发需求通常经历三步：\textbf{codebase exploration}、\textbf{plan}、\textbf{execution}。其中 exploration 负责密集读取代码与上下文，plan 把目标拆成可执行步骤，execution 再根据执行反馈持续推进。
\end{importantbox}

这也是为什么 planning 不只是一个 UI 功能。它实际上是 agent loop 中的状态表示：模型需要把``当前要做什么''、``已经做到哪一步''、``是否需要重新规划''保留在上下文里，并在必要时展示给用户。

\subsection{本章小结}

本讲讨论的是 Codex 里两套互补的规划机制：一种是显式的 Plan Mode，另一种是执行期的 \texttt{update\_plan} 工具。前者偏协作和需求澄清，后者偏运行时进度管理。

